\documentclass[11pt]{article}
\usepackage[margin=1in]{geometry}
\usepackage{enumitem}
% =============================================================================
% Preambles/header.tex — shared LaTeX/Beamer preamble for this project
%
% Use from a Beamer lecture by prepending:
%     \input{header}
% and compiling with TEXINPUTS=../Preambles:$TEXINPUTS (the /compile-latex
% skill does this automatically).
%
% PALETTE CONTRACT. Color names defined here MUST match the names in
% Quarto/theme-template.scss so Beamer and Quarto renderings use the same
% palette. The scripts/check-palette-sync.sh script enforces this.
%
% To customize: edit the HEX values below AND the matching $variable in
% Quarto/theme-template.scss, then run scripts/check-palette-sync.sh.
% =============================================================================

% -----------------------------------------------------------------------------
% Engine detection + encoding
%
% Our /compile-latex skill uses XeLaTeX, where inputenc is unnecessary and
% can produce encoding warnings. Only load inputenc under pdfTeX.
% -----------------------------------------------------------------------------
\usepackage{iftex}
\ifPDFTeX
  \usepackage[utf8]{inputenc}
\fi

\usepackage{amsmath, amssymb, amsthm}
\usepackage{booktabs}
\usepackage{graphicx}

% -----------------------------------------------------------------------------
% PALETTE — mirrors Quarto/theme-template.scss variable names.
% Load xcolor and define colors BEFORE anything that references them
% (notably hyperref, hypersetup, and the Beamer color assignments).
% -----------------------------------------------------------------------------
\usepackage{xcolor}

\definecolor{primary-blue}{HTML}{012169}      % headings, accents
\definecolor{primary-gold}{HTML}{B9975B}      % emphasis, borders
\definecolor{highlight-yellow}{HTML}{F2A900}  % markers, alerts
\definecolor{light-bg}{HTML}{E8EDF5}          % title / section backgrounds
\definecolor{jet}{HTML}{1A1A1A}               % body text

% Semantic colors (used in diagrams and annotations)
\definecolor{positive}{HTML}{15803D}          % good / observed / identified
\definecolor{negative}{HTML}{B91C1C}          % bad / problematic / bias
\definecolor{neutral}{HTML}{525252}           % reference / context

% Highlight accents (match .hi-* classes in the SCSS)
\definecolor{hi-slate}{HTML}{314F4F}
\definecolor{hi-green}{HTML}{2E8B57}
\definecolor{hi-red}{HTML}{C41E3A}

% Semantic aliases so skill templates and snippets can write
% \textcolor{accent}{...} without hard-coding "primary-blue".
\colorlet{accent}{primary-blue}
\colorlet{accent2}{primary-gold}

% -----------------------------------------------------------------------------
% Hyperref — loaded AFTER xcolor + palette so link colors resolve.
% -----------------------------------------------------------------------------
\usepackage{hyperref}
\hypersetup{colorlinks=true, linkcolor=primary-blue, urlcolor=primary-blue,
            citecolor=primary-blue, pdfborder={0 0 0}}

% -----------------------------------------------------------------------------
% Course code — set once per deck (after \input{header}, before \begin{document})
% via \coursecode{CS401}. Shown in the footer on every slide so a deck's course
% is identifiable without opening the title page. Empty by default.
% -----------------------------------------------------------------------------
\makeatletter
\newcommand{\coursecode}[1]{\gdef\@coursecodetext{#1}}
\gdef\@coursecodetext{}
\makeatother

% -----------------------------------------------------------------------------
% Beamer theme assignments (only apply under Beamer).
%
% We detect Beamer via \@ifclassloaded{beamer}, which is the documented,
% non-internal way. This must be wrapped in \makeatletter/\makeatother
% because \@ifclassloaded contains an @-catcode character.
% -----------------------------------------------------------------------------
\makeatletter
\@ifclassloaded{beamer}{%
  \usefonttheme{professionalfonts}
  \setbeamercolor{structure}{fg=primary-blue}
  \setbeamercolor{frametitle}{fg=primary-blue}
  \setbeamercolor{title}{fg=primary-blue}
  \setbeamercolor{subtitle}{fg=primary-gold}
  \setbeamercolor{author}{fg=primary-blue}
  \setbeamercolor{institute}{fg=neutral}
  \setbeamercolor{date}{fg=primary-gold}
  \setbeamercolor{itemize item}{fg=primary-blue}
  \setbeamercolor{itemize subitem}{fg=primary-gold}
  \setbeamercolor{alerted text}{fg=primary-gold}
  \setbeamercolor{block title}{fg=white, bg=primary-blue}
  \setbeamercolor{block body}{bg=light-bg}
  % Semantic block variants: block=definition/concept (blue), exampleblock=
  % worked example (green/positive), alertblock=key takeaway or caution (gold).
  % Keep to <=2 colored boxes per slide across ALL three variants (INV-7).
  \setbeamercolor{block title example}{fg=white, bg=positive}
  \setbeamercolor{block body example}{bg=light-bg}
  \setbeamercolor{block title alerted}{fg=white, bg=primary-gold}
  \setbeamercolor{block body alerted}{bg=light-bg}

  % Minimal footer: course code (left) + page X / N (right), no navigation clutter
  \setbeamertemplate{navigation symbols}{}
  \setbeamertemplate{footline}{%
    \color{neutral}\footnotesize\hspace{0.7em}\@coursecodetext\hfill
    \insertframenumber/\inserttotalframenumber\hspace{0.7em}\vspace{0.3em}}
}{}
\makeatother

% -----------------------------------------------------------------------------
% TikZ setup — libraries the snippet gallery uses
% -----------------------------------------------------------------------------
\usepackage{tikz}
\usetikzlibrary{arrows.meta, positioning, calc, decorations.pathreplacing,
                fit, shapes.geometric, backgrounds}

% Shared TikZ styles — snippets and hand-written diagrams can pick these up
% via `\begin{tikzpicture}[every node/.style={...}]` or by naming specific
% styles (e.g., `\node[dag-node] (X) at (0,0) {X};`).
\tikzset{
  % Node styles for causal DAGs and similar
  dag-node/.style = {
    draw=primary-blue, fill=light-bg, rounded corners=2pt,
    minimum width=1.2cm, minimum height=0.9cm, align=center, font=\small
  },
  decision-node/.style = {
    draw=primary-gold, fill=white, diamond, aspect=1.8,
    minimum width=1.8cm, minimum height=1.2cm, align=center, font=\small,
    inner sep=1pt
  },
  % Edge styles — observed vs counterfactual
  observed-edge/.style = {-{Latex[length=2.5mm]}, thick, draw=jet},
  counterfactual-edge/.style = {-{Latex[length=2.5mm]}, thick, draw=neutral, dashed},
  confound-edge/.style = {-{Latex[length=2.5mm]}, thick, draw=negative, dashed},
  % Data-point styles for plots
  observed-dot/.style = {fill=primary-blue, draw=primary-blue, circle, inner sep=1.2pt},
  counterfactual-dot/.style = {fill=white, draw=neutral, circle, inner sep=1.2pt}
}

% -----------------------------------------------------------------------------
% Convenience macros
% -----------------------------------------------------------------------------
\newcommand{\muted}[1]{\textcolor{neutral}{#1}}
\newcommand{\key}[1]{\textcolor{primary-gold}{\textbf{#1}}}
\newcommand{\good}[1]{\textcolor{positive}{#1}}
\newcommand{\bad}[1]{\textcolor{negative}{#1}}

% Standout frame macro: full-bleed dark background for section transitions.
% Beamer-only (uses \begin{frame}/\setbeamercolor/\usebeamerfont) -- do not
% call from an article-class file (e.g. Notes/<CODE>/*-notes.tex). \muted,
% \key, \good, \bad, and \coursecode above are plain xcolor/gdef and are
% safe to use from either class; verified when Notes/article-class support
% was added.
% Expand: \transitionslide{Your transition title}
\newcommand{\transitionslide}[1]{%
  \begingroup
  \setbeamercolor{background canvas}{bg=primary-blue}
  \begin{frame}[plain]
    \centering
    \vfill
    {\usebeamerfont{title}\color{white}\Large #1\par}
    \vfill
  \end{frame}
  \endgroup
}

% Section divider frame (Beamer-only), an alternative to \transitionslide with a
% darker two-line style: near-black (jet) background, a small white label line
% above a larger gold title, and a thin gold rule along the bottom edge.
% Expand: \sectiondivider{Section 2}{Pointers}
\newcommand{\sectiondivider}[2]{%
  \begingroup
  \setbeamercolor{background canvas}{bg=jet}
  \begin{frame}[plain]
    \vspace*{3.0cm}
    {\color{white}\ttfamily\large #1\par}
    \vspace{0.5cm}
    {\color{primary-gold}\LARGE #2\par}
    \begin{tikzpicture}[remember picture, overlay]
      \draw[primary-gold, line width=2.5pt]
        ($(current page.south west)+(0,0.1)$) -- ($(current page.south east)+(0,0.1)$);
    \end{tikzpicture}
  \end{frame}
  \endgroup
}

% -----------------------------------------------------------------------------
% Channel watermark — "Concepts That Click" (YouTube).
%
% Article-class files (CompetitiveExam Q/A sets, Notes, Assignments) get a
% light diagonal draft watermark on every page plus a centered footer naming
% the channel. Beamer decks get a subtle TikZ background watermark instead
% (the footline already carries the course code). Centralized here so every
% MCQ PDF is branded without per-file edits.
% -----------------------------------------------------------------------------
\makeatletter
\@ifclassloaded{article}{%
  \usepackage{draftwatermark}
  \SetWatermarkText{Concepts That Click}
  \SetWatermarkScale{1}
  \SetWatermarkFontSize{2cm}
  \SetWatermarkLightness{0.86}
  \SetWatermarkAngle{45}
  \usepackage{fancyhdr}
  \pagestyle{fancy}
  \fancyhf{}
  \fancyfoot[C]{\footnotesize\textcolor{neutral}{Concepts That Click~|~YouTube: \href{https://www.youtube.com/@concepts.thatclick}{@concepts.thatclick}}}
  \renewcommand{\headrulewidth}{0pt}
  \renewcommand{\footrulewidth}{0pt}
  \fancypagestyle{plain}{%
    \fancyhf{}%
    \fancyfoot[C]{\footnotesize\textcolor{neutral}{Concepts That Click~|~YouTube: \href{https://www.youtube.com/@concepts.thatclick}{@concepts.thatclick}}}%
    \renewcommand{\headrulewidth}{0pt}%
    \renewcommand{\footrulewidth}{0pt}%
  }
}{}
\@ifclassloaded{beamer}{%
  \setbeamertemplate{background}{%
    \begin{tikzpicture}[remember picture, overlay]
      \node[rotate=45, opacity=0.055, align=center]
        at (current page.center)
        {\Huge\textcolor{neutral}{Concepts That Click}};
    \end{tikzpicture}%
  }
}{}
\makeatother 
\coursecode{JKSSB}
\renewcommand{\thesection}{31.\arabic{section}}
\newtheorem{example}{Example}[section]
\newenvironment{definitionbox}[1]{%
  \par\noindent\textbf{Definition (#1).}\ \itshape
}{\par\vspace{0.3em}}
\title{PowerPoint Basics --- Detailed Lecture Notes}
\author{JKSSB Computer Awareness}
\date{\today}

\begin{document}
\maketitle

\section{What PowerPoint is}
MS PowerPoint is the presentation program of the MS Office suite. A
\textbf{presentation} is a file made of one or more pages, and each page is
called a \textbf{slide}. The slides are shown one after another as a
\textbf{slide show} on a screen or projector. The official syllabus writes
this unit as ``Powerpoint Presentation'', and the usual short form is
\textbf{PPT}. A single slide can hold text, pictures, charts, tables, audio,
and video.

\begin{definitionbox}{Presentation, slide, slide show}
A presentation is the whole file; a slide is one page of it; a slide show is
the full-screen sequential display of those slides to an audience.
\end{definitionbox}

\begin{center}
\begin{tabular}{p{0.24\textwidth}p{0.66\textwidth}}
\toprule
\textbf{Term} & \textbf{Meaning in this lesson} \\
\midrule
Presentation & The whole PowerPoint file, made of slides. \\
Slide & One page of the presentation. \\
Placeholder & The dotted-line box on a slide that holds content. \\
Layout & The arrangement of placeholders on a slide. \\
Theme & A coordinated set of colours, fonts, and effects. \\
Template & A ready-made starter file with a theme, layouts, and sample content. \\
Slide Master & The top-level slide that controls all linked layouts and slides. \\
Notes & The speaker's reminders below a slide, unseen by the audience. \\
Handout & A printed sheet carrying several slides for the audience. \\
\bottomrule
\end{tabular}
\end{center}

The rest of this lesson uses these terms in one consistent sense. A
\textbf{slide} is the page, a \textbf{placeholder} is the box that holds
content on that page, and a \textbf{layout} decides where those boxes sit. A
\textbf{theme} restyles the file, a \textbf{template} starts a new file from a
ready design, and the \textbf{Slide Master} controls the design of every
linked slide at once.

\section{Full forms and file extensions}
Three full forms carry direct-recall marks. \textbf{PPT} stands for PowerPoint
Presentation. \textbf{PPTX} stands for PowerPoint XML Presentation and is the
default format since Office 2007. \textbf{PPSX} stands for PowerPoint Slide
Show. Two more extensions complete the family: \textbf{POTX} is the PowerPoint
Template, and \textbf{PPTM} is the macro-enabled presentation.

\begin{definitionbox}{PowerPoint file extensions}
The editable working file uses \texttt{.pptx}; the template uses
\texttt{.potx}; the show file that opens directly as a slide show uses
\texttt{.ppsx}; the macro-enabled file uses \texttt{.pptm}. Presentations can
also be saved or exported as \texttt{.pdf}.
\end{definitionbox}

The working distinction is simple. Double-clicking a \texttt{.pptx} file opens
it for editing, while double-clicking a \texttt{.ppsx} file launches the
slide show directly. A \texttt{.potx} file is never the finished talk; it is
the reusable pattern from which new talks are started. Confusing these three
extensions is one of the most common wrong answers in this unit.

\section{Slides, placeholders, and layouts}
A slide is a single page of a presentation. Content is not typed freely onto
the page background; it goes into \textbf{placeholders}, the dotted-line boxes
marked Title, Text, Picture, Chart, or Table. Clicking inside a placeholder
allows typing, and the small icons inside a content placeholder insert a
picture, chart, table, or SmartArt graphic directly.

\begin{definitionbox}{Placeholder}
A dotted-line container on a slide that holds one kind of content (title,
text, picture, chart, table, or media).
\end{definitionbox}

A \textbf{layout} is the arrangement of those placeholders. The common layouts
are Title Slide, Title and Content, Section Header, Two Content, Comparison,
Title Only, and Blank. A new slide is added from the Home tab through New
Slide, which also offers the layout gallery; the layout of an existing slide
is changed through Home tab to Layout. The Blank layout carries no
placeholders at all, so every object on it must be inserted by hand.

\begin{example}
A student builds a five-slide class talk on the office desktop. The first
slide uses the Title Slide layout (one title plus one subtitle placeholder);
each following point uses Title and Content; the closing slide reuses Section
Header. Changing the layout of slide 3 from Title and Content to Two Content
rearranges only slide 3; nothing else in the file moves.
\end{example}

\section{Themes, templates, and the Slide Master}
A \textbf{theme} is a coordinated set of colours, fonts, and effects applied
from the Design tab. One click restyles every slide in the file, because the
theme is a design rule, not content. A \textbf{template} is a ready-made
pattern file (extension \texttt{.potx}) that bundles a theme with a set of
layouts and sometimes sample text. Starting a new presentation from a template
reuses that whole design instead of building it again.

\begin{definitionbox}{Theme and template}
A theme is a coordinated set of colours, fonts, and effects. A template is a
reusable starter file (\texttt{.potx}) carrying a theme plus layouts and
sample content.
\end{definitionbox}

The \textbf{Slide Master} is the top-level slide that controls the design of
all layouts and slides linked to it. Editing the font, the logo position, or
the footer on the Slide Master changes every linked slide at once, which is
far faster than editing forty slides one by one. It is opened from the View
tab through Slide Master and left through Close Master View. PowerPoint has
three masters in total: the Slide Master for on-screen slides, the
\textbf{Handout Master} for printed handout layouts, and the \textbf{Notes
Master} for printed notes pages.

\begin{example}
An item states, ``The Handout Master controls the design of all on-screen
slides.'' This is wrong. The Slide Master controls the on-screen slides; the
Handout Master controls only the printed handout layouts. The wrong option has
swapped two of the three masters.
\end{example}

\section{Views: Normal, Outline, Sorter, Reading, Show, Notes Page}
PowerPoint offers several views, each with one job. \textbf{Normal view} is
the default editing view, with three areas: the large slide pane in the
centre, the slides or outline pane with thumbnails on the left, and the notes
pane below. All typing, formatting, and inserting happen here. \textbf{Outline
View} shows only the text outline of every slide, which suits restructuring
long talks.

\textbf{Slide Sorter view} shows thumbnails of all slides together, so slides
can be reordered by drag and drop and checked, deleted, or duplicated at a
glance. \textbf{Reading View} shows the show almost full-screen for reading on
the author's own screen, with simple navigation but without editing tools.
\textbf{Slide Show view} is the real full-screen presentation to an audience.
\textbf{Notes Page view} shows one slide above its speaker-notes text, one
page per slide, exactly as it will print.

\begin{definitionbox}{Normal view}
The default editing view of PowerPoint, combining the slide pane, the
slides or outline pane, and the notes pane.
\end{definitionbox}

\begin{example}
A student has twelve slides in the wrong order. The fastest fix is Slide
Sorter view: drag the thumbnails into the right sequence. Opening each slide
in Normal view and cutting and pasting would work but wastes time; choosing
Reading View would not help at all, because Reading View cannot reorder
anything.
\end{example}

\section{Notes, handouts, and printing}
\textbf{Notes} (speaker notes) are reminders typed in the notes pane below a
slide. The audience never sees them during the show; they exist so the speaker
does not forget a figure, a definition, or the next point. During the show the
speaker reads them only in Presenter View on the speaker's own screen. Notes
print through File to Print to Notes Pages.

\begin{definitionbox}{Speaker notes and handouts}
Speaker notes help the speaker and stay hidden from the audience. Handouts
help the audience and are printed sheets carrying several slides per page.
\end{definitionbox}

\textbf{Handouts} are those printed sheets: 1, 2, 3, 4, 6, or 9 slides per
page. The most examined choice is 6 slides per page, and the 3-per-page layout
adds ruled lines beside each slide so the audience can write. Their design is
controlled by the Handout Master. The full print path is File to Print, with
choices for Full Page Slides, Handouts, Notes Pages, and Outline View.

\section{Running and navigating the slide show}
Two shortcuts start the show. \textbf{F5} begins from the first slide, and
\textbf{Shift+F5} begins from the current slide. This pair is a confirmed past
item (WO 2026 Q80 tested F5 for ``from the beginning''). During the show, N,
Space, Right arrow, or Enter moves forward; P or Left arrow moves back; Esc
ends the show. Pressing B blanks the screen to black and pressing W blanks it
to white, which is useful while the speaker talks without a slide. The Slide
Show tab holds timing, rehearsal, and monitor settings.

A \textbf{transition} is the movement between one slide and the next (Fade,
Push, Wipe, Morph), set on the Transitions tab with one transition per slide.
An \textbf{animation} is the movement of an object on a slide, set on the
Animations tab: Entrance, Emphasis, Exit, and Motion Paths, managed in the
Animation Pane. SmartArt conversion behaviour (which content cannot convert
without re-entry) has also been examined (JA 2026 Q77), so candidates should
practise the Insert to SmartArt path on text and picture lists rather than
assume every object converts.

\section{Exam retrieval and common confusions}
Six distinctions prevent most errors on this topic.
\begin{enumerate}
  \item The \textbf{slide} is the page; the \textbf{placeholder} is the
    dotted box holding content; the \textbf{layout} arranges those boxes.
  \item \textbf{.pptx} is the editable file, \textbf{.ppsx} is the show file,
    and \textbf{.potx} is the template. Do not swap them.
  \item A \textbf{theme} restyles (colours, fonts, effects); a
    \textbf{template} is the reusable starter file.
  \item The \textbf{Slide Master} controls on-screen slides; the Handout and
    Notes Masters control print layouts.
  \item \textbf{Normal} is the default editing view; \textbf{Slide Sorter}
    reorders; Reading reviews; \textbf{notes} serve the speaker and
    \textbf{handouts} serve the audience.
  \item \textbf{F5} starts from the beginning and \textbf{Shift+F5} from the
    current slide; a \textbf{transition} moves between slides while an
    \textbf{animation} moves objects on a slide.
\end{enumerate}

\begin{example}
An item asks which view is the default editing view and offers Normal, Slide
Sorter, Reading, and Slide Show. Trace the definitions: Normal is the editing
view with the slide, thumbnail, and notes panes; Slide Sorter manages order;
Reading reviews; Slide Show presents. The answer is Normal view. If the stem
instead asks which view reorders slides by drag and drop, the answer changes
to Slide Sorter.
\end{example}

\subsection*{Exercises}
\begin{enumerate}[label=\arabic*.]
  \item Distinguish a slide, a placeholder, and a layout, giving one example
    of each.
  \item Expand PPT, PPTX, and PPSX, and state the file extension of a
    PowerPoint template.
  \item Distinguish a theme from a template.
  \item What does the Slide Master control, and how do the Handout and Notes
    Masters differ from it?
  \item List four PowerPoint views and state the job of each.
  \item Distinguish speaker notes from handouts, and state which viewers each
    serves.
  \item Which keys start a show from the beginning and from the current slide,
    and which key ends the show?
  \item Distinguish a transition from an animation.
\end{enumerate}

\subsection*{Solutions}
\begingroup\small
\begin{enumerate}[label=\arabic*.]
  \item A slide is one page of the presentation (for example, slide 2 of a
    class talk). A placeholder is the dotted-line box holding content on that
    page (for example, the Title placeholder). A layout is the arrangement of
    those boxes (for example, Title and Content).
  \item PPT is PowerPoint Presentation; PPTX is PowerPoint XML Presentation;
    PPSX is PowerPoint Slide Show. A PowerPoint template uses
    \texttt{.potx}.
  \item A theme is a coordinated set of colours, fonts, and effects that
    restyles a file; a template is a reusable starter file
    (\texttt{.potx}) that bundles a theme with layouts and sample content.
  \item The Slide Master controls the design of all linked on-screen slides
    and layouts. The Handout Master controls printed handout layouts and the
    Notes Master controls printed notes pages; neither restyles the on-screen
    slides.
  \item Normal is the default editing view; Slide Sorter shows thumbnails for
    reordering; Reading View shows the show for on-screen review; Slide Show
    presents full-screen; Notes Page shows one slide above its notes text.
    Any four with correct jobs earn full credit.
  \item Speaker notes are the speaker's hidden reminders below a slide; they
    serve the speaker. Handouts are printed multi-slide sheets; they serve
    the audience.
  \item F5 starts from the beginning; Shift+F5 starts from the current slide;
    Esc ends the show.
  \item A transition is the movement between one slide and the next; an
    animation is the movement of an object on a slide (Entrance, Emphasis,
    Exit, Motion Paths).
\end{enumerate}
\endgroup
\end{document}
